\documentclass[10pt,a4paper]{amsart}
\usepackage[margin=19mm]{geometry}
\usepackage{amsmath,amssymb,mathtools}
\usepackage[scr=rsfs,cal=euler]{mathalfa}
\usepackage{xfrac}
\usepackage[unicode,hidelinks,bookmarks=false]{hyperref}
\newcommand{\ie}{i.e.\ }
\newcommand{\cf}{cf.\ }
\newcommand{\resp}{respectively\ }
\newcommand{\Subsection}{\subsection}
\providecommand{\nobreakdash}{\nobreak\hbox{-}}
\setlength{\emergencystretch}{2em}
\multlinegap0pt
\usepackage{bm}
\usepackage{algorithm}\usepackage{algpseudocode}
\newtheorem{lemma}{Lemma}
\title{RTD-Lite: Scalable Topological Analysis for Comparing Weighted Graphs in Learning Tasks}
\author{Eduard Tulchinskii, Daria Voronkova, Ilya Trofimov, Evgeny Burnaev, Serguei Barannikov}
\date{Source: arXiv:2503.11910v1 (CC BY 4.0)}
\begin{document}
\maketitle

\noindent\emph{Source and licence.} RTD-Lite: Scalable Topological Analysis for Comparing Weighted Graphs in Learning Tasks. Version arXiv:2503.11910v1 (CC BY 4.0), distributed by arXiv under the Creative Commons Attribution 4.0 International licence, http://creativecommons.org/licenses/by/4.0/, as recorded in the arXiv licence metadata for this exact version and shown on the abstract page https://arxiv.org/abs/2503.11910v1. Full licence notice: \texttt{licenses/arxiv-2503.11910-CC-BY-4.0.txt}.\par\smallskip

\noindent\textit{Editorial scope.} Counted unit: the article's lemma lemma (source0.tex lines 1387--1389). The article's complete original proof of it is delivered with it (source0.tex lines 1391--1434, one \texttt{proof} environment of 44 lines). No mathematics is written, repaired or re-derived: the statement and the proof are the article's own lines, in the article's own notation, and no retained line is substituted or re-flowed.  No further source-local context is needed: every cross-reference of every retained line resolves inside the two delivered spans.  Nothing was omitted from any retained span, so the package carries no omission record.  No retained line contains a citation, so the wrapper carries no bibliography and no bibliography entry.  No retained line contains a figure environment, an image-inclusion command or drawing code, so no image is delivered and the package declares no asset dependency.  Editorial formatting only, in addition: an AMS \texttt{amsart} preamble that restates the article's own declarations and macros used by the retained lines, namely the environments \texttt{lemma} on separate counters, together with the standard packages those lines need.  The retained environments print in this document as Lemma 1 in order of appearance, whereas the article prints the same items with the numbering of its own full document.  The complete native source of the article as distributed in the arXiv e-print for this exact version is bundled unchanged as \texttt{sources/174800/source0.tex}.
\begin{lemma}
Dimensionality of $\ker{(r_{0, \alpha})}$ equals to the number of intervals in RTD-Lite barcode that contains \nolinebreak $\alpha$.
\end{lemma}

\begin{proof}




By the rank-nullity theorem for the map 

\[
r_{0,\alpha}: H_0(A^{\leq\alpha}) \rightarrow H_0(C^{\leq\alpha}),
\]

we obtain: 

\begin{equation}
    \dim{H_0(A^{\leq\alpha})} = \dim{\ker{(r_{0, \alpha})}} + \dim{\mathrm{Im}(r_0,\alpha)}.
    \label{app:eq1}
\end{equation}

By construction of the graph $C$, $r_{0,\alpha}$ is epimorphism, i.e. $\mathrm{Im}(r_0,\alpha) = H_0(C^{\leq\alpha})$. Thus, Equation \ref{app:eq1} is equivalent to:

\begin{equation}
    \dim{\ker{(r_{0, \alpha})}} = \dim{H_0(A^{\leq\alpha})} - \dim{H_0(C^{\leq\alpha})}.
    \label{app:eq2}
\end{equation}

Let us see how $\dim{\ker{(r_{0, \alpha})}}$ changes when increasing the value of $\alpha$ via analysing the dimensionalities of homology groups $H_0(A^{\leq\alpha}), H_0(C^{\leq\alpha})$. 
First, note that for $\alpha = 0$, both $A^{\leq\alpha}$ and $C^{\leq\alpha}$ consist only of $n$ vertices with no edges. Thus, $\dim{H_0(A^{\leq\alpha})} = \dim{H_0(C^{\leq\alpha})} = n$ and $\dim{\ker{(r_{0, \alpha})}} = 0$.
Also, the intersection of the RTD-Lite barcode with $\alpha=0$ line is empty.

Now let's add to the set of vertices the edges of $C\supseteq A$  one by one in the increasing weight order.  Let us increase the value of $\alpha$ from $\alpha_1$ to the next weight $\alpha_2$ of an edge from $C$, $\alpha_2 > \alpha_1$. First, consider the case when the new edge $e$ appends to $A^{\leq\alpha_2}$ compared to $A^{\leq\alpha_1}$. If $e$ connects the two vertices from the same connected component, then $\dim{H_0(A^{\leq\alpha_2})} = \dim{H_0(A^{\leq\alpha_1})}$ and $\dim{H_0(C^{\leq\alpha_2})} = \dim{H_0(C^{\leq\alpha_1})}$ and $\dim{\ker{(r_{0, \alpha_2})}} = \dim{\ker{(r_{0, \alpha_1})}}$. If $e$ connects the vertices from the two different connected components, then the two components merge together and $\dim{H_0(A^{\leq\alpha_2})} = \dim{H_0(A^{\leq\alpha_1})} - 1$. If $e$ corresponds to the vertices from the same connected component in $C^{\leq\alpha_1}$, then $\dim{H_0(C^{\leq\alpha_2})} = \dim{H_0(C^{\leq\alpha_1})}$. In this case, $\dim{\ker{(r_{0, \alpha_2})}} = \dim{H_0(A^{\leq\alpha_2})} - \dim{H_0(C^{\leq\alpha_2})} = (\dim{H_0(A^{\leq\alpha_1})} - 1) - \dim{H_0(C^{\leq\alpha_1})} = \dim{\ker{(r_{0, \alpha_1})}} - 1$. If $e$ corresponds to the vertices from different connected components, then $\dim{H_0(C^{\leq\alpha_2})} = \dim{H_0(C^{\leq\alpha_1})} - 1$ and $\dim{\ker{(r_{0, \alpha_2})}} = \dim{H_0(A^{\leq\alpha_2})} - \dim{H_0(C^{\leq\alpha_2})} = (\dim{H_0(A^{\leq\alpha_1})} - 1) - (\dim{H_0(C^{\leq\alpha_1})} - 1) = \dim{\ker{(r_{0, \alpha_1})}}$. 

Second, consider the case when $A^{\leq\alpha_2} = A^{\leq\alpha_1}$ (and $H_0(A^{\leq\alpha_2}) = H_0(A^{\leq\alpha_1})$) but a new edge $e$ appears in $C^{\leq\alpha_2}$ compared to $C^{\leq\alpha_1}$. If $e$ corresponds to the vertices from the same connected component, then $H_0(C^{\leq\alpha_2}) = H_0(C^{\leq\alpha_1})$ and $\dim{\ker{(r_{0, \alpha_2})}} = \dim{\ker{(r_{0, \alpha_1})}}$. If $e$ corresponds to the vertices from different connected components, then $H_0(C^{\leq\alpha_2}) = H_0(C^{\leq\alpha_1}) - 1$ and $\dim{\ker{(r_{0, \alpha_2})}} = \dim{H_0(A^{\leq\alpha_2})} - \dim{H_0(C^{\leq\alpha_2})} = 
\dim{H_0(A^{\leq\alpha_1})} - (\dim{H_0(C^{\leq\alpha_1})} - 1) = \dim{\ker{(r_{0, \alpha_1})}} + 1$. 

Thus, there are two cases of a change in $\dim{\ker{(r_{0, \alpha_1})}}$:
\begin{itemize}
    \item $\dim{\ker{(r_{0, \alpha_2})}} = \dim{\ker{(r_{0, \alpha_1})}} - 1$.
    In this case, the two connected components that were merged in $C^{\leq\alpha}$ for some $\alpha < \alpha_2$ now merge in $A^{\leq\alpha_2}$. This corresponds to closing of the interval $(\alpha, \alpha_2)$ in RTD-Lite barcode.
    \item $\dim{\ker{(r_{0, \alpha_2})}} = \dim{\ker{(r_{0, \alpha_1})}} + 1$.
    In this case, there are two connected components that merge in $C^{\leq\alpha_2}$ and not (yet) merge in $A^{\leq\alpha_2}$. This corresponds to opening of the interval $(\alpha_2, \alpha)$ in RTD-Lite barcode for some $\alpha > \alpha_2$.
\end{itemize}

Therefore, the decrease and increase in \(\dim \ker (r_0, \alpha)\) correspond to the closing and opening of intervals in the RTD-Lite barcode. Thus, for each \(\alpha\), \(\dim \ker (r_0, \alpha)\) equals the number of intervals in the RTD-Lite barcode that contain \(\alpha\).
\end{proof}

\end{document}
